\documentclass[12pt,a4paper]{article}

\usepackage[utf8]{inputenc}
\usepackage[T1]{fontenc}
\usepackage[portuguese]{babel}
\usepackage[left=3cm,top=3cm,right=2cm,bottom=2cm]{geometry}
\usepackage{graphicx}                   % inserir imagem
\usepackage{listings}                   % formatar bloco de código
\usepackage{helvet}                     % fonte sans serif
\usepackage{hyperref}
\hypersetup{
    colorlinks=true,
    urlcolor=blue,
}

\begin{document}
\fontfamily{phv}\selectfont

\section{Nome do Projeto}
PlaceZ (PLEI-CÊS)

\section{Componentes do Time}
\begin{itemize}
    \item Guilherme Amaral Giffoni
    \item Pedro Henrique Ferreira
    \item Rodrigo Fernandes Moraes
\end{itemize}

\section{\href{https://github.com/guigiffoni/PlaceZ}{Link para o Github}}

\section{Proposta de Desenvolvimento}
\begin{itemize}
    % o que é e para que serve
    \item Objetivo do App \\
    Facilitar a organização de eventos com muitas pessoas, como um encontro
    entre amigos para um restaurante/lanchonete, uma festa de aniversário,
    RPGs...

    % quem irá se beneficiar com a utilização do App
    \item Público alvo do App  \\
    Grupos de amigos, organizadores de festas, restaurantes, bares, lanchonetes,
    e correlatos.

    \item Principais funcionalidades \\
    Encontrar lugares a partir de proximidade com a localização do grupo;
    filtrar por: preço, tipo de comida e restrição alimentar; organizar a partir
    de um ``calendário compartilhado'' o melhor dia e horário para o evento com 
    cada membro.

    A filtragem pode ser tão específica quanto necessário.

    \item Rascunho das telas do App \\
    Não temos
    
\end{itemize}

\end{document}